\section{Conclusion}
\label{sec:conclusion}

This paper establishes the partisan neutrality of the \ar{} algorithm through
ensemble comparison.  The main findings are:

\begin{enumerate}
  \item \ar{} is not systematically biased toward either party.  It produces
    outcomes within one seat of the ensemble median in five of six studied
    states.  The median percentile of Democratic seats across all states is
    the 53rd percentile---essentially indistinguishable from the geographic
    median.

  \item The Georgia deviation (38th percentile of Democratic seats) reflects
    geographic sorting, not algorithmic bias.  The Rodden concentration
    effect reduces the number of compact Democratic districts in any compact
    algorithm applied to Georgia.

  \item The \ar{} plan falls inside the proportionality corridor for the
    competitive states studied (NC, WI, PA, MN).  This is a strong result:
    an algorithm with no proportionality objective nevertheless produces
    proportional-adjacent outcomes in competitive states.

  \item Algorithm choice matters in sorted states.  Short-burst algorithms
    with compactness-correlated objectives can produce extreme partisan
    outcomes.  \ar{}'s edge-cut objective is not correlated with partisan
    outcomes in the way that some compactness metrics are.
\end{enumerate}

The practical implication: the question ``is this map partisan?'' can be
answered quantitatively for \ar{} plans using ensemble comparison.  For
competitive states, the answer is clearly no.  For sorted states like Georgia,
the answer is more nuanced: the plan produces a Republican-leaning outcome
relative to proportional representation, but this lean is consistent with
the geographic constraints and is present in a majority of valid plans.

Courts and legislators can use this evidence to conclude that the \ar{}
plan's partisan outcomes are not a product of manipulation, even in states
where those outcomes favor one party.
